\section{Rò rỉ khái niệm}
\label{sec:ch15-ro-ri}

\index{rò rỉ!hiện tượng quan sát được|(}%
Mục trước khép lại trên một câu hỏi mà phép tối ưu được phép trả lời: đặt giá
trị nào vào $\hat c_k$ để đầu phân loại đoán đúng $y$. Mục này đọc câu trả lời,
và đọc nó trước ở chỗ nhìn thấy được chứ không ở chỗ định nghĩa.

\index{rò rỉ!phân bố giá trị kích hoạt}%
Cảnh nhìn thấy được là phân bố các giá trị $\hat c_1$ trên tập kiểm
tra~\autocite{p30leakage}. Bài của mục~\ref{sec:ch15-do-ro-ri} lấy một nhiệm vụ
đúng hai khái niệm
nhị phân, nhãn là $y = c_1$ hoặc $c_2$, rồi huấn luyện ba mô hình trên đó: một
mô hình cứng, và hai mô hình mềm với $\beta$ đặt ở mức trung bình rồi ở mức rất
thấp. Với mô hình cứng, phân bố của $\hat c_1$ có hai cụm nằm gần 0 và gần 1,
tức gần đúng phân bố nhị phân của $c_1$. Với mô hình mềm ở mức $\beta$ trung
bình, hai cụm ấy tách tiếp: bên trong cụm gần 1 xuất hiện thêm cấu trúc, và cấu
trúc ấy phân biệt được theo giá trị của khái niệm còn lại và theo nhãn. Với mô
hình mềm ở mức $\beta$ rất thấp, phân bố của $\hat c_1$ giống phân bố của $y$
hơn là giống phân bố của $c_1$.

Một đại lượng được huấn luyện để nhận khái niệm
thứ nhất, và trên tập kiểm tra nó phân bố như nhãn nhiệm vụ. Người giám sát nhìn
vào cột ấy, thấy một số trong khoảng từ 0 đến 1, và đọc nó như \enquote{mô hình
tin khái niệm thứ nhất có mặt tới mức này}. Cách đọc ấy sai, không phải vì con
số vô nghĩa mà vì nó mang hai thứ chồng lên nhau và người đọc chỉ được kể về
một.

\index{rò rỉ!định nghĩa}%
Hiện tượng ấy có tên. \tn{Rò rỉ}{leakage} xảy ra khi mô hình cất thêm thông tin vào các giá
trị khái niệm học được, ngoài chính trạng thái của khái niệm, và đầu phân loại
dùng phần thông tin thêm ấy để đoán nhãn. Nó có lợi cho phép tối ưu, vì mất mát
trên nhãn giảm mà mất mát trên khái niệm không tăng bao nhiêu, nên phép tối ưu
tìm ra nó mà không cần ai bảo. Cái nó phá là điều duy nhất mà cả họ mô hình này
được xây để bảo đảm: rằng đọc bảng $\hat c_k$ là đọc được lý do mô hình ra quyết
định. Chữ ấy đã xuất hiện một lần trong sách với nghĩa khác, ở
\tn{rò rỉ tương tác}{interaction leakage} của
chương~\ref{ch:tro-choi-attribution}, nơi cái rò đi từ một hiệu ứng tương tác
sang một \tn{hiệu ứng riêng}{pure individual effect} bên trong một bảng điểm số.
Hai hiện tượng khác nhau
mang cùng một chữ tiếng Anh, và trong chương này thì \enquote{rò rỉ} viết không
kèm bổ ngữ luôn chỉ hiện tượng vừa mô tả.

\index{rò rỉ!hai loại}%
Bài ấy chia rò rỉ thành hai loại theo chỗ thông tin thêm ấy nói về cái
gì. \tn{Rò rỉ khái niệm-nhiệm vụ}{concepts-task leakage} là khi $\hat c_k$ mang
thêm thông tin về nhãn $y$;
đây là loại chính, vì nó là loại trực tiếp giúp mô hình đoán nhiệm vụ tốt hơn.
\tn{Rò rỉ liên khái niệm}{interconcept leakage} là khi $\hat c_k$ mang thêm thông tin về những khái niệm
khác. Loại thứ hai thường phụ hơn, và nó không hoàn toàn vô ích cho mô hình:
biết $\hat c_2$ từ $\hat c_1$ cho mô hình một đường vòng để vẫn đoán đúng nhiệm
vụ khi một khái niệm bị nhận sai.

\index{rò rỉ!hiện tượng quan sát được|)}%
Hai chữ \enquote{thêm} trong hai câu vừa rồi cần một mốc so sánh, và mốc ấy là
điều làm cả khung này chạy được. Trong một bộ dữ liệu có gán nhãn khái niệm,
lượng thông tin mà khái niệm thật $c_k$ mang về nhãn thật $y$ là một đại lượng
tính được, không phụ thuộc vào mô hình nào. Nó là mốc. Nếu giá trị học được
$\hat c_k$ mang nhiều thông tin về $y$ hơn mốc ấy, phần chênh là thông tin mà mô
hình đã tự thêm vào, vì dữ liệu không có sẵn để mà thêm. Mục~\ref{sec:ch15-do-ro-ri}
biến câu ấy thành hai điểm số.

\index{rò rỉ!can thiệp mất chỗ dựa}%
Hình~\ref{fig:ch15-nut-that-va-ro-ri} vẽ cùng một chuyện ở chỗ nó xảy ra, và vẽ
kèm hệ quả mà mục~\ref{sec:ch15-dung-cu-duoc-kiem} sẽ dùng làm dụng cụ đo. Một
đơn vị của tầng thắt mang hai phần: phần trạng thái khái niệm, thứ người giám
sát đọc được, và phần thông tin thêm, thứ người giám sát không đọc được vì không
có cột nào cho nó. Đầu phân loại đọc cả hai. Khi người giám sát can thiệp, tức
đặt $\hat c_k$ về giá trị thật $c_k$, thì phần dưới biến mất cùng phần trên, vì
chúng là cùng một con số. Đầu phân loại vốn đã học cách trông cậy vào phần dưới
nay không còn nó, nên độ chính xác nhiệm vụ có thể tụt đi sau một phép can thiệp
đúng. Đó là dấu hiệu thô nhất của rò rỉ và cũng là dấu hiệu dễ hiểu nhất: sửa
cho mô hình một chỗ sai mà nó lại đoán tệ hơn.

\begin{figure}[htbp]
  \centering
  \begin{tikzpicture}[
  every node/.style={font=\footnotesize, align=center},
  box/.style={draw=black!65, rounded corners=2pt, fill=black!5,
    minimum width=24mm, minimum height=12mm},
  shown/.style={draw=black!85, line width=1pt, fill=black!12,
    minimum width=30mm, minimum height=9mm},
  hidden/.style={draw=black!55, dotted, thick, fill=white,
    minimum width=30mm, minimum height=9mm},
  flow/.style={-{Stealth[length=2mm]}, thick, black!70},
  thinflow/.style={-{Stealth[length=1.6mm]}, black!55},
  lbl/.style={font=\scriptsize, align=center, black!55}
]
  \node[box] (enc) at (0,0) {Bộ mã hóa\\khái niệm};

  \node[shown]  (top) at (4.6,0.45)  {trạng thái khái niệm $k$};
  \node[hidden] (bot) at (4.6,-0.45) {thông tin thêm về $y$};
  \node[draw=black!85, line width=1pt, rounded corners=2pt,
        fit=(top)(bot), inner sep=1.6mm] (cell) {};
  \node[lbl, left=1.5mm of cell, text width=13mm] {một số\\$\hat c_k$};

  \node[box] (head) at (9.4,0) {Đầu\\phân loại};
  \node[box, minimum width=14mm] (out) at (12.3,0) {$\hat y$};

  \node[lbl, left=2mm of enc] {$x$};

  \draw[flow] (enc.east) -- (cell.west);
  \draw[flow] (cell.east) -- (head.west);
  \draw[flow] (head) -- (out);

  \node[lbl, above=1.5mm of top, text width=44mm]
    {người giám sát đọc được chừng này};
  \node[lbl, below=1.5mm of bot, text width=44mm]
    {và không đọc được chừng này};

  \draw[thinflow] (2.2,-2.3) -- (2.2,0.45) -- (top.west);
  \node[lbl] at (2.2,-2.7) {can thiệp};
  \node[lbl] at (7.4,-2.7) {đặt $\hat c_k$ về giá trị thật $c_k$;\\ô dưới đi theo, vì nó là cùng một số};
\end{tikzpicture}

  \caption{Rò rỉ không đi vòng qua tầng thắt mà đi xuyên qua nó. Khung ngoài là
  một đơn vị của tầng thắt, tức đúng một con số $\hat c_k$, và hai ô bên trong
  là hai phần thông tin mà con số ấy chở chứ không phải hai đường đi: người giám
  sát đọc được ô trên và không có cột nào cho ô dưới, còn đầu phân loại nhận cả
  khung. Bảng khái niệm vì thế vẫn đầy đủ về hình thức trong khi đã thiếu về nội
  dung, và không phép kiểm nào trên bảng ấy phát hiện được. Ô dưới vẽ nét chấm
  theo quy ước của sách cho thứ vắng mặt khỏi cái được trưng ra.}
  \label{fig:ch15-nut-that-va-ro-ri}
\end{figure}

\index{ví dụ chạy suốt sách!ngưỡng hóa thành nhiệm vụ nhị phân}%
Nhiệm vụ hai khái niệm ở trên tính được bằng tay, và nó chính là hàm mà sách đã
dùng từ chương~\ref{ch:giai-thich-de-lam-gi}. Hàm ấy là
$f(x_1,x_2) = x_1 + 2x_2 + x_1x_2$ trên hai \tn{đặc trưng}{feature} bật tắt,
nhận bốn giá trị
$0$, $1$, $2$ và $4$ tại bốn tổ hợp. Lấy ngưỡng tại $1$, nghĩa là đặt $y = 1$
khi $f \ge 1$, thì $y$ bằng 0 đúng ở tổ hợp $(0,0)$ và bằng 1 ở ba tổ hợp còn
lại: đó là phép hoặc của hai đặc trưng, tức đúng nhiệm vụ mà bài báo dùng làm ví
dụ. Hai đặc trưng của hàm chạy suốt sách trở thành hai khái niệm, và ngưỡng biến
một hàm số thực thành một nhiệm vụ phân loại nhị phân.

Còn thiếu một giả thiết nữa mới tính được, và nó là giả thiết của sách chứ không
của bài báo: lấy $c_1$ và $c_2$ độc lập, mỗi cái bằng 0 hoặc 1 với xác suất như
nhau. Bộ dữ liệu tổng hợp của bài báo sinh hai khái niệm có tương quan, nên các
con số dưới đây minh họa cách đo chứ không phải là con số bài báo báo cáo.

Với giả thiết ấy thì $y = 0$ chỉ ở một trong bốn tổ hợp, nên $y$ bằng 1 với xác
suất $3/4$. Entropy của nhãn là
\[
  H(y) = -\tfrac{3}{4}\log_2\tfrac{3}{4} - \tfrac{1}{4}\log_2\tfrac{1}{4}
       = 0.811 \text{ bit} .
\]
Ở đây $H$ là entropy, đo lượng bất định của một biến ngẫu nhiên tính bằng bit.
Biết $c_1 = 1$ thì $y$ chắc chắn bằng 1, nên bất định còn lại bằng 0; biết
$c_1 = 0$ thì $y$ bằng $c_2$, nên bất định còn lại bằng đúng 1 bit. Trung bình
hai trường hợp cho $0.5$ bit, và lượng thông tin mà $c_1$ mang về $y$ là phần
bất định bị xóa đi,
\[
  I(c_1, y) = H(y) - 0.5 = 0.311 \text{ bit} .
\]
Ký hiệu $I$ với hai đối số là thông tin tương hỗ, đại lượng đo phần bất định của
biến này mà biết biến kia xóa được. Chương~\ref{ch:do-do-trung-thuc} dùng chữ
$I$ trần cho một \tn{phép nhiễu}{perturbation} trong định nghĩa infidelity; hai
nghĩa tách nhau bằng hình dạng, vì thông tin tương hỗ luôn đi kèm hai đối số còn
phép nhiễu kia luôn đứng một mình.

Chia cho $H(y)$ để con số không phụ thuộc thang đo của nhãn, mốc của khái niệm
thứ nhất trên nhiệm vụ này là
\[
  \frac{I(c_1, y)}{H(y)} = \frac{0.311}{0.811} = 0.384 .
\]
Một mô hình cứng học đúng khái niệm cho $\hat c_1 = c_1$, nên tỉ số của nó cũng
là $0.384$ và phần chênh bằng 0. Ở đầu kia, một giá trị học được đã thành bản
sao hoàn hảo của nhãn mang trọn $H(y)$, nên tỉ số bằng 1 và phần chênh là
$1 - 0.384 = 0.616$. Vậy trên nhiệm vụ này, phần chênh chạy trong khoảng từ 0
tới $0.616$, và biết một mô hình đứng ở đâu trong khoảng ấy là biết bao nhiêu
phần thông tin trong cột $\hat c_1$ không nói về khái niệm thứ nhất. Với hai
khái niệm độc lập thì $I(c_1,c_2) = 0$, nên mốc liên khái niệm bằng 0 và
\emph{mọi} mức phụ thuộc học được giữa $\hat c_1$ và $\hat c_2$ đều là rò rỉ.

\index{độ trung thực!rò rỉ là bản sao trong họ mô hình nào}%
Bây giờ mới nói được câu mà nhan đề chương gợi ra, và nói kèm chỗ nó dừng. Đọc
theo bảng~\ref{tab:ch01-bon-quan-he}, rò rỉ đúng là họ hàng của
\tn{độ trung thực}{faithfulness}
theo định nghĩa~\ref{def:ch01-do-trung-thuc}: cả hai đều so cái được trưng ra
với cái mà mô hình thật sự dựa vào, và cả hai đều hỏng theo cùng một kiểu, là
cái được trưng ra thiếu chứ không sai. Nhưng ba chỗ khác nhau đủ lớn để không
được gộp. Thứ nhất, ở đây cái được trưng ra không phải một lời giải thích do một
phương pháp bên ngoài dựng lên mà là một tầng của chính mô hình, nên không có
tập can thiệp nào phải công bố theo yêu cầu mà
mục~\ref{sec:ch14-dieu-chua-giai-quyet} để lại. Thứ hai, rò rỉ là một đại lượng
định nghĩa so với một tập khái niệm đã gán nhãn: đổi tập khái niệm thì mốc đổi
theo, nên nó nói \enquote{cột này mang thêm bao nhiêu so với tập ấy} chứ không
nói \enquote{mô hình thật sự dựa vào cái gì}. Thứ ba, không bài nào trong ba bài
của chương gọi nó là độ trung thực, và chữ \enquote{faithful} không xuất hiện
lần nào trong phần thân của bài nào. Câu ở đầu mục này là câu của sách, và nó
đứng được đúng tới đây.
